\part{Bonus Lectures}

\section{Inverse and Direct Limits}

\subsection{Inverse Limits}

\begin{definition}[\textbf{Inverse System}]
	Fix a category \( \cC \), we view \( (I, \leq) \) as an ordered set, and viewed as a Category. An
	inverse system in \( \cC \) is a \textbf{functor}
	\[
		(I,\leq) \to \cC^{\text{op}}
	\]
	for some ordered set \( (I,\leq) \). Explicitly, this is given by \( (x_i)_{i\in I} \) with \( x_i
	\in \ob(\cC)\) + for every \( i\leq j \), a map
	\[
		\begin{aligned}
			\vp_{ij}: x_j & \to x_i                 \\
			\vp_{ij}      & \in \Hom_{\cC}(x_j,x_i)
		\end{aligned}
	\]
	s.t. it is compatible with the morphisms
	\[
		\begin{aligned}
			 &
			\begin{cases}
				\vp_{ii} = \Id \; \forall \; i \\
				\forall \; i \leq j \leq k
			\end{cases}                                \\
			 & x_k \xrightarrow{\vp_{jk}}  x_j \xrightarrow{\vp_{ij}} x_i \\
			 & \vp_{ik} = \vp_{ij}\circ \vp_{jk}
		\end{aligned}
	\]
\end{definition}

\begin{definition}[\textbf{Inverse Limits}]
	The \underline{inverse limit} (projective limit), denoted as \( \vpl_{i\in I} X_i \) of such an
	inverse/projective system is given by
	\[
		\vpl X_i \xrightarrow{p_i} X_i \quad \forall \; i \in I
	\]
	that satisfies: \( i\leq j \implies \) the following diagram being commutative:
	\[
		\begin{tikzcd}
			\vpl X_i \arrow[r, "p_i"] \arrow[d, "p_j"'] & X_i \\
			X_j \arrow[ru, "\vp_{ij}"'] &
		\end{tikzcd}
	\]
	and it is universal with this property: \( \forall \; Y \in \ob(\cC) \) with maps \( Y
	\xrightarrow{q_i} X_i \) that satisfy \( \vp_{ij}\circ q_j = q_i \; \forall \; i \leq j \), there
	exists a unique morphism \( f \), such that the following diagram is commutative
	\[
		\begin{tikzcd}
			\vpl X_i \arrow[r, "f"] \arrow[d, "p_i"'] & Y \arrow[ld, "q_i"] \\
			X_i &
		\end{tikzcd}
	\]
	namely \( q_i \circ f = p_i \; \forall \; i\).
\end{definition}

\begin{eg}
	\leavevmode
	\begin{enumerate}
		\item Any \( I \), with \( \leq \) being \( = \), we have
		      \[
			      (I,=) \to \cC^\text{op}
		      \]
		      is just a family \( (X_i)_{i\in I} \) and \( \vpl_{i\in I}X_i = \prod_{i\in I} X_i \), which
		      is exactly the direct product.
		\item Let \( (\ZZ_{\geq 0}, \leq) \) and \( \cC = \{\text{commutative ring}\} \), then let \( R \) be any
		      ring, we have an inverse system as follows:
		      \[
			      \cdots \twoheadrightarrow \underbrace{\qo{R[x]}{(x^{n+1})}}_{=S_n} \twoheadrightarrow \cdots
			      \twoheadrightarrow \underbrace{\qo{R[x]}{(x^2)}}_{=S_1} \twoheadrightarrow \underbrace{\qo{R[x]}{(x)}}_{=S_0}
		      \]
		      the inverse limits in the category of the commutative rings.
	\end{enumerate}
\end{eg}

\begin{remark}
	Suppose have the inverse system given by:
	\[
		\begin{aligned}
			(I,\leq) & \to \text{\textit{Rings}} \\
			i        & \to R_i                   \\
			i\leq j  & \to \vp_{ij}: R_j \to R_i
		\end{aligned}
	\]
	with when \( j\leq k \) we have the following diagram being commutative:
	\[
		\begin{tikzcd}
			\vpl R_i \arrow[r] \arrow[d] & R_k \arrow[ld] \\
			R_j &
		\end{tikzcd}
	\]
	then
	\[
		\vpl_{i\in I} R \coloneqq  \{a = (a_i)_i \in \prod_{i\in I} R_i \; | \; a_i = \vp_{ij}(a_{ij})
		\; \forall \; i \leq j \} \subseteq \prod_{i\in I} R_i \xrightarrow{\pi_i} R_i
	\]
\end{remark}

\begin{exercise}
	\leavevmode
	\begin{enumerate}
		\item \( \vpl R_i \subseteq \prod_{i\in I} R_i \) and being a subring.
		\item The maps \( p_i = \pi_i \circ \text{incl} \) satisfy:
		      \[
			      \forall \; i \leq j, p_i = \vp_{ij} \circ p_j
		      \]
		\item This is universal with these properties: follows from the universal property of \(
		      \prod_{i\in I} R_i \).
	\end{enumerate}
\end{exercise}

\begin{eg}
	\leavevmode
	\[
		\vpl_{n\geq 0} \frac{R[x]}{(x^{n+1})}
	\]
	and
	\[
		\qo{R[x]}{(x^{n+1})} \to \qo{R[x]}{(x^n)} \cong \; \R[[x]]
	\]
	where the maps:
	\[
		R[[x]] \to \qo{R[x]}{(x^n)} \cong \qo{R[[x]]}{(x^n)}
	\]
\end{eg}

\begin{remark}
	The same construction via direct products works in the category of \underline{Sets}, \( R
	\)-\underline{mod}, etc.
\end{remark}

\subsection{Direct Limits}

\begin{definition}[\textbf{Direct System}]
	A \underline{direct system} in a category \( \cC \) is a functor \( I \xrightarrow{F} \cC \). Explicitly:
	\[
		\begin{aligned}
			i \in I & \rsa X_i \in \ob(\cC)                               \\
			i\leq j & \rsa \vp_{ij}: X_i \to X_j \in \Hom_{\cC}(X_i, X_j)
		\end{aligned}
	\]
	s.t.
	\begin{itemize}
		\item \( \vp_{ij} = \Id \; \forall \; i \).
		\item \( i\leq j \leq k \implies \vp_{jk} \circ \vp_{ij} = \vp_{ik} \).
	\end{itemize}
\end{definition}

\begin{definition}[\textbf{Direct Limit}]
	The \underline{direct limit} (injective limit), denoted as \( \varinjlim_{i \in I} X_i \) is the
	dual notion of that of inverse limit: namely \( \forall \; i \), s.t. for \( i\leq j \), the following diagram is commutative:
	\[
		\begin{tikzcd}
			X_i \arrow[r, "f_i"] \arrow[d, "\vp_{ij}"] & \varinjlim_{i\in I} X_i \\
			X_j \arrow[ru, "f_j"] &
		\end{tikzcd}
	\]
	and this is universal: \( \forall \; Y \) with \( X_i \xrightarrow{g_i} Y \) that satisfy the same
	compatible condition,
	\[
		\exists \; ! \; u_i: \varinjlim X_i \to Y \; \text{ s.t. } u\circ f_i = g_i \quad \forall \; i
	\]
\end{definition}

\begin{eg}
	If \( \cC = \) \( R \)-\underline{mod}, and \( ``\leq''\) is \( ``='' \), then we have seen that
	this is \( \bigoplus_{i\in I} M_i \). Direct limits are easier to handle and have better
	properties if \( (I, \leq) \) is \underline{filtering}: \( \forall \; i,j \in I, \; \exists\; k
	\in I \), s.t. \( i,j \leq k \).
	\begin{eg}
		If \( A \) be any set, then
		\[
			(I = \{B \subseteq A \; | \; B \text{ finite. }\}, \subseteq) \text{ is filtering.}
		\]
	\end{eg}
\end{eg}

\begin{proposition}
	Filtering direct limits exist in the category of rings.
\end{proposition}

\begin{proof}
	Given a direct system \( (R_i)_{i\in I} \) and \( \vp_{ij} : R_i \to R_j \) for \( i \leq j \), we
	define
	\[
		\varinjlim_{i\in I} R_i \coloneqq  \bigsqcup_{i\in I} \qo{R_i}{\sim}
	\]
	where \( a\in R_{j_1}\sim b \in R_{j_2} \)  iff \( \exists\; k \), s.t. \( j_1, j_2 \leq k \) and
	\[
		\vp_{j_1 k}(a) = \vp_{j_{2}k} (a)
	\]
	We can define \( (+) \) and \( (\cdot) \) as follows: if \( a\in R_{j_1} \) and \( b\in R_{j_2}
	\), then
	\[
		\overline{a} + \overline{b} \coloneqq \vp_{j_1 k} (a) + \vp_{j_2 k} (b)
	\]
	where \( k\) is s.t. \( j_1, j_2 \leq k \), and
	\[
		\overline{a} \cdot \overline{b} \coloneqq \vp_{j_1 k} (a) \vp_{j_2 k} (b)
	\]
	\begin{exercise}
		\leavevmode
		\begin{enumerate}
			\item Check the well-definedness and this make \( \varinjlim_{i\in I} R_i \) a ring.
			\item This satisfy the universal property.
		\end{enumerate}
	\end{exercise}
	We have the induce maps:
	\[
		\begin{tikzcd}
			R_j \arrow[r] \arrow[rd, "f_j"] & \bigsqcup_{i\in I} R_i \arrow[d]\\
			& \varinjlim R_i
		\end{tikzcd}
	\]
	when is \( f_j(a) = 0 \)? see that \( f_j (a) = 0 \iff \exists \; k, j \leq k \), s.t. \( \vp_{jk} (a) = 0 \). In particular,
	if all \( \vp_{jk} \) are injective \( \implies  \) all \( f_j \) are injective \( \implies
	\varinjlim_{i\in I} R_i = \bigcup_{i\in I} R_i \).
\end{proof}

\begin{eg}
	Fix commutative ring \( R \), suppose \( A \) is a set, and
	\[
		I = \{B \subseteq A \; | \; B \text{ finite.}\}
	\]
	For every \( B \in I, B \subseteq B' \) in \( I \rsa \)
	\[
		R[x_b \; ; \; b \in B] \hookrightarrow R[x_b \; ; \; b \in B']
	\]
	this is a direct system and \( I \) is filtering. Have
	\[
		R[x_b \; ; \; b\in A] \coloneqq \varinjlim_{\substack{B \subseteq A \\ \text{ finite }}} R [x_b
				\; ; \; b \in B] = ``\bigcup_{\substack{B \subseteq A \\ \text{ finite }}}'' R[x_b \; ; \; b \in
				B]
	\]
	Have the universal property of \( R[x_b \; ; \; b \in A] \): \( \forall \) commutative \( R
	\)-algebra \( S \), and every map \( f: A \to S \), \( \exists \;! R \)-algebra homomorphism
	\[
		\begin{aligned}
			u : R[x_b \; ; \; b \in A] & \to S                           \\
			\text{s.t. } u(x_b)        & = f(b) \quad \forall \; b \in A
		\end{aligned}
	\]
	this follows similar property for finite \( A \) + universal property of \( \varinjlim \).
\end{eg}

\begin{remark}
	In the category of \( R \)-\underline{mod} where \( R \) is any ring, one can use a similar
	construction. There is another one that works for \underline{all} \( I \). If \( (M_i)_{i\in I}
	\), and \( M_i \xrightarrow{\vp_{ij}} M_j \) is a direct system of \( R \)-\underline{mod}, then take
	\[
		\begin{tikzcd}
			\bigoplus_{i\in I} M_i \arrow[r] & \qo{\bigoplus_{i\in I} M_i}{\angl{\alpha_k (\vp_{jk}(a)) -
					\alpha_j (a) \; | \; j \leq k, a\in M_j}} \\
			M_j \arrow[u, "\alpha_j"] \arrow[ru, "f_j"]&
		\end{tikzcd}
	\]
	it is easy to check: this satisfy the definition of \( \varinjlim_{i\in I} M_i \).
\end{remark}

\begin{theorem}
	For every field \( K \), \( \exists \) algebraic extension \( K \hookrightarrow \overline{K} \),
	s.t. \( \overline{K} \) is algebraically closed.
\end{theorem}

\begin{proof}
	\leavevmode
	\begin{itemize}
		\item Step 1: construct an algebraic extension \( K \hookrightarrow K_1 \), s.t. every \( f\in
		      K[x] \) of degree \( \geq 1 \) has a root in \( K_1 \), define
		      \[
			      A = \{f\in K[x] \; | \; \deg(f) \geq 1\}
		      \]
		      Take \( R = K[y_f \; | \; f\in A] \) and take \( J = (f_{(y_f)} \; | \; f \in A) \)
		      \begin{claim}
			      \( J \neq R \).
		      \end{claim}
		      Suppose \( J = R \implies \)
		      \[
			      \begin{aligned}
				      \exists \; f_1                       & , \ldots, f_r \in A \\
				      g_1                                  & , \ldots, g_r \in A \\
				      \text{ s.t. } \sum g_i f_i (y_{f_i}) & = 1
			      \end{aligned}
		      \]
		      we say: there exists \( K \hookrightarrow L_1 \hookrightarrow L_2 \cdots L_r \), s.t. \( f_1 \) has a root in \( L_1 \),
		      \( f_2 \) has a root in \( L_2 \), so \( \exists \; a_1,\ldots, a_r \in L_r \), s.t. \(
		      f_i (a_i) = 0 \; \forall \; i \). By the universal property of \( R \), \( \exists \)
		      morphism of \( K \)-algebra
		      \[
			      \begin{aligned}
				      R = K[y_f \; | \; f \in A]     & \to L_r \\
				      y_{f_i}                        & \to a_i \\
				      g \ne f_1, \ldots, f_r, \; y_g & \to 0
			      \end{aligned}
		      \]
		      and we apply this to the above equation, and in \( L_r \) we shall have \( 1 = 0 \),
		      leading to a contradiction \( \lightning \). Thus we have the conclusion.
		      \begin{conclusion}
			      \( J = (f_{(y_f)} \; | \; f \in A) \ne R  \implies \exists\) maximal ideal \( M \) in \( R \), s.t.
			      \( J \subseteq M \). Thus
			      \[
				      \begin{aligned}
					      K                 & \hookrightarrow K_1 = \qo{R}{M} \text{ field.} \\
					      f(\overline{y_f}) & = 0 \text{ in } K_1
				      \end{aligned}
			      \]
		      \end{conclusion}

		      \begin{note}
			      \( \qo{K_1}{K} \) is algebraic. By the following:
			      \[
				      \begin{cases}
					      K_1 = K(\overline{y_f} \; ; \; f \in A) \\
					      \text{Each } \overline{y_f} \text{ is algebraic over } K
				      \end{cases}
				      \implies \qo{K_1}{K} \text{ is algebraic.}
			      \]
		      \end{note}
		\item Step 2: Repeat this construction to get field extension:
		      \[
			      K_0 = K \hookrightarrow K_1 \hookrightarrow K_2 \hookrightarrow \cdots
		      \]
		      s.t.
		      \begin{itemize}
			      \item \( \qo{K_{n+1}}{K_n} \) is algebraic \( \forall \; n \geq 0 \).
			      \item \( \forall \; n \geq 0, \forall \; f \in K_n[X] \) of degree \( \geq 1 \), \( f \)
			            has a root in \( K_{n+1} \).
		      \end{itemize}
		      Take \( \overline{K} = \varinjlim_{i\geq 0} K_i \), have \( K_i \hookrightarrow
		      \overline{K} \), s.t. \( \overline{K} = \bigcup_{i \geq 0} K_i \).
		      \begin{itemize}
			      \item \( \overline{K} \) is a field since each \( K_i \) is a field and \( \overline{K} = \bigcup_{i\geq 0} K_i \).
			      \item \( \qo{\overline{K}}{K} \) is \underline{algebraic}: \( \forall \; a \in \overline
			            K\), \( \exists \; i \), s.t. \( a\in K_i \) and \( \qo{K_i}{K} \) is algebraic by
			            induction on \( i \) since \( \qo{K_j}{K_{j-1}} \) is algebraic.
			      \item \( \overline K \) is \underline{algebraically closed}: \( \forall \; f \in
			            \overline K[x] \), \( \exists \; i \geq 0 \), s.t. \( f\in K_i [x]  \implies f\) has a
			            root in \( K_{i+1} \subseteq \overline K \).
		      \end{itemize}
	\end{itemize}
\end{proof}

